% paper/sections/03_method.tex
%
% First draft. Structure follows paper/03_method.md (six subsections:
% three-arm tool surface, harness design, evaluation integrity, capability
% overlap framing, cache + cost methodology, artifact-suite contract).
%
% Drafting rules:
%   * No bare digits in prose — paper/lint.py enforces this; numbers come
%     through \result / \respct / \resp macros or live in §4/§5.
%   * Arm names are consolidated paper convention (baseline / bash_only /
%     code_only) per the 2026-05-27 decision in paper/outline.md L257.
%     Legacy harness names (tool_rich / onlycode) are confined to §4.2.
%   * §3.3 omits the planned audit paragraph: no CSV produces
%     \result{audit_total} / \result{audit_propagated} and §5's
%     restructured outline has no audit subsection to forward into.
%   * Labels match those already referenced in sections/01_introduction.tex
%     (sec:method-integrity, sec:method-artifact).

% ─────────────────────── 3.1 Three-arm tool surface ────────────────────
\subsection{Three-arm tool surface}
\label{sec:method-arms}

Each arm is defined by what tools are exposed to the agent at process
launch; the same Claude Code or Codex CLI binary runs all three. The
\texttt{baseline} arm passes no restriction flag and exposes the agent's
default surface --- Read, Grep, Glob, Edit, Write, Bash, and the
agent-built-in subagents on Claude Code; the analogous file
primitives, shell tool, and structured \texttt{apply\_patch} on Codex.
The \texttt{bash\_only} arm restricts to bash plus read-only browsing
primitives, disallowing every file-editing built-in
(\texttt{Edit}/\texttt{Write}/\texttt{MultiEdit}/\texttt{NotebookEdit}
on Claude; \texttt{shell\_tool} with structured patching disabled on
Codex). The \texttt{code\_only} arm replaces the entire surface with a
single MCP tool, \texttt{mcp\_\_codebox\_\_execute\_code}, served by
our exec-server stack as a persistent Python+Bash REPL keyed by
working directory across calls; all native built-ins are explicitly
disallowed.

Two implementation asymmetries are worth surfacing. First, Codex
exposes no feature flag for disabling its structured
\texttt{apply\_patch} tool, so on Codex the \texttt{bash\_only} and
\texttt{code\_only} arms enforce restriction with a prompt-prefix
directive rather than a hard flag; this is soft enforcement and model
compliance is not guaranteed. We measure and report the resulting
leakage in §\ref{sec:results} and find it empirically
negligible. Second, \texttt{code\_only} depends on REPL state
persisting between \texttt{execute\_code} invocations: without a
persistent kernel, the round-trip savings of ``one script per task''
collapse, so the exec-server pins
the persistent-kernel environment flag for the runs that feed §\ref{sec:results}.

The \texttt{bash\_only} arm matches the tool surface of
mini-SWE-agent~\cite{repo:miniswagent} but does not reuse its loop,
prompt template, or scaffold; we apply the bash-only restriction to
the same Claude Code and Codex binaries the other two arms use, so the
contrast isolates the tool surface variable rather than comparing
scaffold implementations.

% ──────────────────────── 3.2 Harness design ───────────────────────────
\subsection{Harness design}
\label{sec:method-harness}

\paragraph{Per-instance setup (SWE-bench).} Each SWE-bench instance is
cloned at its \texttt{base\_commit} and a per-instance Python virtual
environment is built against it. The repository and the virtual
environment are each given their own pristine fuse-overlayfs lower
directory, so a per-arm upper directory captures every write and is
discarded between arms --- no re-clone, no
\texttt{pip~install~-e~.} between arms. The
virtual-environment overlay is mounted back at its original creation
path because console-script shebangs bake that absolute path at
install time. Artifact tasks are self-contained and skip this step.

\paragraph{Git-history strip.} Before the agent runs, the worktree is
collapsed to a single orphan commit at its current tree; all refs,
packed refs, reflogs, and alternates are deleted, and
\texttt{git~gc~--prune=now} repacks the resulting object set. The
agent cannot recover the upstream reference fix via \texttt{git~log},
\texttt{git~show}, or the reflog. Fixed author and committer dates
make the orphan SHA deterministic for a given tree, so repeated runs
strip to the same state.

\paragraph{Overlay refresh between arms.} \texttt{git~reset} cannot
un-create files an earlier arm added (fuse-overlayfs copy-up plus
EEXIST semantics), so the harness unmounts, deletes the upper and
work directories, recreates them, remounts, and re-strips history
between arms. The agent's view of the merged mount path is stable;
the underlying state is fresh.

\paragraph{Per-arm subprocess isolation.} Each agent invocation
creates a fresh temporary configuration directory containing only the
credentials and minimal config needed for the run, and points the
binary at it via \texttt{CLAUDE\_CONFIG\_DIR} or
\texttt{CODEX\_HOME}. No session-persistence files, no cross-run
state. Claude is invoked with
\texttt{--dangerously-skip-permissions\,--no-session-persistence};
Codex with \texttt{--dangerously-bypass-approvals-and-sandbox}.

\paragraph{Wall-time cap.} Each invocation runs in a fresh process
group; on timeout the harness sends \texttt{SIGKILL} to the entire
group (so subagents and spawned shells die with the parent) and the
run is recorded \texttt{FAIL} with a synthetic \texttt{wall\_timeout}
JSONL record. The numerical wall budget is pinned in
§\ref{sec:setup}.

% ────────────────── 3.3 Evaluation integrity (SWE-bench) ───────────────
\subsection{Evaluation integrity}
\label{sec:method-integrity}
\label{sec:m-integrity}

The canonical SWE-bench protocol applies the gold \texttt{test\_patch}
only after the agent terminates: at run time the agent observes the
repository at \texttt{base\_commit}, and the tests used for grading
are not present in the worktree~\cite{jimenez2024swebench}. Our
harness implements deferred application explicitly. The
\texttt{test\_patch} is applied only after the agent's process exits,
and a \texttt{pytest~--collect-only} gate then runs against the
patched tree to confirm that the held-out tests are collectible. A
zero-item collection scores \texttt{FAIL} rather than letting the test
runner silently report success against a non-existent suite.

Because our agents --- unlike the bash-only patch-submission scaffold
used in the original SWE-bench evaluation --- can write arbitrary
files anywhere in the worktree, the canonical ``apply the patch''
step is preceded by two cleanup operations. Every file that the
\texttt{test\_patch} targets is restored from \texttt{HEAD} to discard
any agent edits to it, and every path the patch lists as a new file
is removed from the worktree so an agent-created file cannot collide
with a held-out test and block the apply. \texttt{git~apply} is then
invoked against the resulting tree. If the apply still fails ---
almost always because an agent-created file could not safely be
removed in the second step --- the instance is recorded
\texttt{FAIL} with reason \emph{post-agent~git~apply~failed:~agent
edits conflicted with the held-out test patch}, rather than scored
against a tree we cannot trust. This is strictly more conservative
than the canonical protocol, which assumes a clean apply.

% ──────────────────── 3.4 Capability Overlap framing ───────────────────
\subsection{The Capability Overlap framing}
\label{sec:method-capability}

We adopt the tool-use tax framing of
Zhang~et~al.~\cite{zhang2026tooltax}: a tool delivers net benefit iff
the task-specific capability gain it provides exceeds the per-call
cost of carrying its definition in the prompt. The original framework
targets general math and question-answering tool use; we apply it to
the coding-agent IDE surface. Most Claude Code IDE primitives are
bash-subsets in \emph{capability} (Read, Grep, Glob, and Write reduce
to \texttt{cat}, \texttt{grep}, \texttt{find}, and shell-heredoc
respectively); \texttt{Edit} is the one primitive with non-overlapping
capability --- atomic byte-precise replacement with linting --- which
is awkward to express in bash. The per-call tax (system-message
bytes spent defining each tool) is task-invariant; the \emph{gain}
term depends on the regime, on which primitives the agent actually
uses, and on the agent's call pattern.

§\ref{sec:results} reads the four-cell cost structure as the
empirical landing of this inequality across (regime, agent) pairs;
§\ref{sec:discussion} decomposes the structure into the three
mechanisms (edit-friction path-cost, failure-cost on doomed
trajectories, and pricing-asymmetry batching/pagination) that the
framework alone does not pin down.

% ──────────────── 3.5 Cache isolation and cost reporting ───────────────
\subsection{Cache isolation and cost reporting}
\label{sec:m-cost}
\label{sec:method-cost}

Across the four surfaces this paper touches --- the Anthropic
Messages API, the Claude Code CLI, the OpenAI Responses API, and the
OpenAI Codex CLI --- no documented parameter, flag, environment
variable, or namespace knob forces a server-side cache miss or scopes
cached tokens to a single caller. The only documented mechanism is
implicit: change the cached content so the cache key changes.
Anthropic's \texttt{cache\_control} controls breakpoints, not
bypass~\cite{docs:anthropic2026promptcaching}; the context-editing
primitive exposes lifecycle and scope but no per-session isolation
knob~\cite{docs:anthropic2026contextediting}; the Claude Code CLI's
\texttt{--no-session-persistence} flag controls local session state,
not the upstream
cache~\cite{docs:claudecode2026cli,docs:claudecode2026settings}.
OpenAI's \texttt{prompt\_cache\_key} is a routing and partition hint,
not a disable
switch~\cite{docs:openai2026promptcaching}, and the Responses API's
stored-state primitives (\texttt{store=true} retention,
\texttt{previous\_response\_id}, the Sessions API) are alternative
stateful channels rather than isolation knobs layered on top of the
cache~\cite{docs:openai2026convstate}; the Codex CLI exposes no cache
flag or environment variable~\cite{repo:openai2026codex}. We
therefore cannot guarantee per-task cache isolation in the reported
numbers and instead address the resulting non-stationarity in the
cost-reporting methodology.

Cost in §\ref{sec:results} is reported as a single
\emph{cache-adjusted} column, derived from the underlying
token-billed formula
$\text{cost} = \text{input\_tokens} \cdot r_{\text{in}}
+ \text{output\_tokens} \cdot r_{\text{out}}$
(every token billed at non-cached rates; the per-model
rates~$r_{\text{in}}, r_{\text{out}}$ are pinned in
§\ref{sec:setup}) plus a per-arm median-floor adjustment.
For each arm, let $M_{\text{arm}}$ be the median across that arm's
instances of first-turn \texttt{cache\_read\_input\_tokens} ---
empirically the floor representing ``system prompt plus tool
definitions are warm in cache.'' Any instance whose first-turn
\texttt{cache\_read} falls below $M_{\text{arm}}$ is bumped up to
$M_{\text{arm}}$ for accounting; instances already at or above the
floor and all multi-turn cache reads are unchanged. The adjustment
normalises the system-prompt warm/cold lottery without claiming we
achieved isolation.

The adjustment deliberately stops at the arm-median, which
corresponds to the size of the shared system prompt plus tool
definitions; we do not claim it controls for cache sharing
\emph{beyond} that prefix. Such post-prompt cross-task sharing is,
however, vanishingly unlikely. Both providers' prompt caches key on
a prefix hash invalidated by any byte difference in the cached span,
and past the system prompt the conversation diverges per task: the
first user message is the task statement (a fresh string per
instance), the model's response is task-conditional, and every
subsequent tool-call payload carries task-specific arguments. For a
cache hit to extend further, two instances would need byte-identical
task statements, model responses, and tool-call payloads in
identical order. Empirically, per-instance first-turn
\texttt{cache\_read} clusters tightly around each arm's system-prompt
token count, with no bimodal tail consistent with cross-task content
reuse. We deliberately do not report raw API-billed cost: it is
sensitive to run order and ambient cache state we cannot control, so
reporting it would mislead readers about effects we cannot
disentangle from confounds.

% ──────────────────── 3.6 The artifact-suite contract ──────────────────
\subsection{The artifact-suite contract}
\label{sec:method-artifact}
\label{sec:m-artifact}

The artifact suite is a new benchmark; this subsection is the
methodology layer that defends it the way §\ref{sec:method-integrity}
defends our SWE-bench protocol. Per-category counts and what each
category probes are reported in §\ref{sec:setup}; this section
specifies the contract, not the inventory.

\paragraph{Task contract.} Every artifact instance is a four-part
bundle on disk: a \texttt{task.yaml} (problem statement, execution
budget fields, and pointers to the other three parts), a
\texttt{workspace/} directory of public files copied into the agent's
scratch dir at materialise time, a \texttt{grader/hidden.py} exposing
a single \texttt{grade(scratch\_dir) $\rightarrow$ GradeResult} entry
point, and a \texttt{grader/reference\_output.*} reference artifact
used to verify the grader itself. Instance IDs follow the convention
\texttt{<category>\_\_<slug>}. The execution-budget fields
(\texttt{max\_code\_runs}, \texttt{max\_wall\_seconds}) are declared
in \texttt{task.yaml} for schema completeness but enforcement is off
in the corpus that feeds §\ref{sec:results}; the only wall-clock cap
is the harness-level budget from §\ref{sec:method-harness}.

\paragraph{Grader contract.} Each task's \texttt{grade(scratch\_dir)}
returns a score in $[0, 1]$ together with a pass flag. Three rules
constrain what \texttt{grade()} may do: it must be
\emph{deterministic} (re-grading the same scratch directory returns
the same score); \emph{offline} (no network calls, no API keys, no
clock-dependent oracles); and \emph{seeded-random only} (any
stochasticity is gated by a fixed seed declared in the grader).
\texttt{grade()} may not write to the scratch directory, which would
contaminate later re-grading. A reference-output round-trip ---
running the grader against \texttt{workspace/} merged with
\texttt{reference\_output/} and asserting that the score equals one
--- is the grader's own self-test and is executed by CI on every
push.

\paragraph{Subprocess isolation.} Graders run in a fresh subprocess
launched per task, with results serialised as JSON on stdout. A
grader-side exception cannot kill the harness, and stdout pollution
from grader code is forbidden by the same contract. This is the
artifact-side analogue of the
\texttt{pytest~--collect-only} pre-flight gate in
§\ref{sec:method-integrity}.

\paragraph{No-leak invariant.} Before the agent runs, the
materialiser copies \texttt{workspace/} into the scratch dir and
scans the resulting tree for any path matching
\texttt{grader/hidden.py} or \texttt{reference\_output.*}. A match
raises and the run is aborted \emph{before} the agent starts; no run
can complete with the answer key in its working directory. The check
is pre-flight by design --- catching a leak post-hoc would already
have biased the grade.
